\SourceLeaf{42}
\[
-gz+\tfrac12c\int\mathrm{d}s.\varphi(\mu,\mathrm{d}s)+C\int\mathrm{d}S.\Phi(\mu,\mathrm{d}S)
\]
ubi $z$ est altitudo puncti $\mu$ supra planum $H$, atque integrale primum per totum spatium $s$, secundum per totum spatium $S$ extendendum est. Qui valor, a solo loco puncti $\mu$ pendens, si per $[\mu]$ denotatur. erit
\[
\Omega=c\int\mathrm{d}s.[\mathrm{d}s]
\]
integratione per totum spatium $s$ extensa.

Brevius hoc ita exprimitur:
\[
\Omega=-gc\int z\mathrm{d}s+\tfrac12cc\int\int\mathrm{d}s.\mathrm{d}s'.\varphi(\mathrm{d}s,\mathrm{d}s')+cC\int\int\mathrm{d}s.\mathrm{d}S.\Phi(\mathrm{d}s,\mathrm{d}S)
\]
ubi $s,s'$ proprie denotant unum idemque spatium (a corpore mobili expletum), sed bis in elementa sua pro duplici integratione resolvendum.

\Article{4}
Corporum fluidorum indoles characteristica consistit in perfecta mobilitate vel minimarum partium, ita ut figuram quamlibet induere possint, et vel minimae potentiae, figuram mutare nitenti, cedant. In fluidis inexpansibilibus (liquidis), quibus nostra disquisitio dicata est, volumen cuiusvis particulae constans manere debet pro omnibus figurae mutationibus. Considerando itaque corpus fluidum, cuius motus per corpus immobile solidum (vas) limitatur, et in cuius particulas praeter gravitatem agere supponimus tum attractionem partium mutuam, tum attractionem partium vasis, status aequilibrii poscet, ut valor ipsius $\Omega$ sit maximum, i. e. ut nulla transpositio infinite parva partium fluidi ipsi $\Omega$ incrementum positivum inducere possit. Quapropter quum manifesto valor ipsius $\Omega$ eatenus tantum mutari possit, quatenus figura spatii, quod totum fluidum implet, mutatur (neque vero per solum motum fluidi internum), aequilibrium aderit, quoties $\Omega$ pro nulla illius figurae mutatione infinite parva cum figura vasis conciliabili, manente volumine constante, augmentum capere potest. Sponte hinc sequitur, si figura omnino nullam mutationem assumere possit (vase fluidum undique cingente et tangente), vires illas in fluidum agentes motum internum fluidi producere non posse, sed sibi aequilibrium facere.

\Article{5}
Progredimur ad accuratiorem investigationem expressionis $\Omega$, quae tamquam fundamentum theoriae aequilibrii fluidorum considerari debet. Incipiendo \SourceLeaf{43}a termino primo, sponte patet, $\int z\mathrm{d}s$ exhibere productum e volumine spatii $s$ in altitudinem centri gravitatis eius supra planum $H$, adeoque $c\int z\mathrm{d}s$ productum massae, $gc\int z\mathrm{d}s$ productum ponderis fluidi in eandem altitudinem. Quodsi itaque partes fluidi praeter gravitatem alii vi non essent obnoxiae, altitudo centri gravitatis in statu aequilibrii esse deberet quam minima, unde facile colligitur, superficiei partem liberam, seu partes liberas, in uno eodemque plano horizontali esse debere, fluidum superne limitante.

\Article{6}
Evolutio termini secundi et tertii refertur ad duos casus particulares problematis generalis, ubi, propositis duobus spatiis quibuscunque, singula elementa primi spatii cum singulis elementis secundi combinari, et producta e ternis factoribus, puta e volumine elementi spatii primi, volumine elementi spatii secundi, et functione data distantiae mutuae, in summam colligi debent. Terminus secundus refertur ad casum eum, ubi ambo spatia identica sunt, tertius ad eum, ubi alterum spatium totum est extra alterum: problema completum duos alios casus complectitur, scilicet ubi vel alterum spatium est pars alterius, vel alterum cum altero partem communem habet. Quamquam vero tum duo priores casus ad institutum nostrum sufficere, tum duo reliqui ad illos facile reduci possent, tamen operae pretium erit, problema per se satis insigne generalitate completa amplecti. Spatia in hac disquisitione generali per $s,S$, functionem distantiae per characteristicam $\varphi$ denotabimus, ita ut in applicatione ad terminum secundum loco ipsius $S$ ipsum spatium $s$, in applicatione ad terminum tertium loco functionis $\varphi$ ipsam $\Phi$ substituere oporteat. Agitur itaque de integrali
\[
\int\int\mathrm{d}s.\mathrm{d}S.\varphi(\mathrm{d}s,\mathrm{d}S)
\]
quod speciem quidem prae se fert integrationis duplicis, sed revera, quum utriusque spatii elementa a ternis variabilibus pendeant, integrationem sextuplicem implicat, quam iam ad integrationem quadruplicem reducere docebimus.

\Article{7}
Initium facimus ab evolutione integralis $\int\mathrm{d}s.\varphi(\mu,\mathrm{d}s)$ per omnes partes spatii $s$ extendendi, denotante $\mu$ punctum determinatum vel extra vel intra spatium $s$ situm. Concipiatur superficies sphaerica radio $=1$ circa centrum $\mu$ de\SourceLeaf{44}scripta, atque in elementa infinite parva divisa; sit $\mathrm{d}\Pi$ tale elementum, secetque recta a $\mu$ versus punctum huius elementi ducta superficiem spatii $s$ deinceps in punctis $p',p'',p'''$ etc., quorum multitudo par erit vel impar, prout $\mu$ extra vel intra spatium $s$ iacet; distantias $\mu p',\mu p'',\mu p'''$ etc. denotabimus per $r',r'',r'''$ etc. Ducantur porro rectae a $\mu$ versus singula puncta peripheriae elementi $\mathrm{d}\Pi$, quo pacto formabitur spatium pyramidale, atque exsecabuntur e superficie spatii $s$, apud puncta $p',p'',p'''$ etc., elementa resp. per $\mathrm{d}t',\mathrm{d}t'',\mathrm{d}t'''$ etc. denotanda. Denique sit $q'$ angulus inter rectam $p'\mu$ atque normalem in elementum $\mathrm{d}t'$ extrorsum ductam; et perinde sint $q'',q'''$ etc. inclinationes similium normalium apud puncta $p'',p'''$ etc. ad rectam versus $\mu$ ductam. Ita manifesto erit
\[
\mathrm{d}\Pi=\pm\frac{\mathrm{d}t'.\cos q'}{r'r'}=\mp\frac{\mathrm{d}t''.\cos q''}{r''r''}=\pm\frac{\mathrm{d}t'''.\cos q'''}{r'''r'''}\ \text{etc.}
\]
valentibus signis superioribus vel inferioribus, prout $\mu$ est extra vel intra spatium $s$.

Porro patet, integrale $\int\mathrm{d}s.\varphi(\mu,\mathrm{d}s)$ pro spatii $s$ partibus intra spatium illud pyramidale contentis haberi per integrale $\mathrm{d}\Pi.\int rr\varphi r.\mathrm{d}r$ extensum ab $r=r'$ usque ad $r=r''$, dein ab $r=r'''$ usque ad $r=r''''$ etc., si $\mu$ iaceat extra spatium $s$, vel extensum ab $r=0$ usque ad $r=r'$, dein ab $r=r''$ usque ad $r=r'''$ etc., si $\mu$ iaceat intra spatium $s$. Quodsi itaque statuimus indefinite
\[
\int rr\varphi r.\mathrm{d}r=-\psi r
\]
constante integrationis ad lubitum accepta, integrale $\int\mathrm{d}s.\varphi(\mu,\mathrm{d}s)$, quatenus extenditur ad partes spatii $s$ intra spatium illud pyramidale sitas, erit
\[
\begin{aligned}
&=\mathrm{d}\Pi.(\psi r'-\psi r''+\psi r'''-\ \text{etc.})\\
&=\frac{\mathrm{d}t'.\cos q'.\psi r'}{r'r'}+\frac{\mathrm{d}t''.\cos q''.\psi r''}{r''r''}+\frac{\mathrm{d}t'''.\cos q'''.\psi r'''}{r'''r'''}+\ \text{etc.}
\end{aligned}
\]
quoties $\mu$ iacet extra spatium $s$; sed
\[
\begin{aligned}
&=\mathrm{d}\Pi.(\psi0-\psi r'+\psi r''-\psi r'''+\ \text{etc.})\\
&=\mathrm{d}\Pi.\psi0+\frac{\mathrm{d}t'.\cos q'.\psi r'}{r'r'}+\frac{\mathrm{d}t''.\cos q''.\psi r''}{r''r''}+\frac{\mathrm{d}t'''.\cos q'''.\psi r'''}{r'''r'''}+\ \text{etc.}
\end{aligned}
\]
quoties $\mu$ iacet intra spatium $s$.

Iam si haec summatio per omnes superficiei sphaericae partes colligitur, integrale $\int\mathrm{d}s.\varphi(\mu,\mathrm{d}s)$ completum fit
\SourceLeaf{45}
\[
\begin{aligned}
\text{in casu priori}\qquad&=\int\frac{\mathrm{d}t.\cos q.\psi r}{rr}\\
\text{in casu posteriori}\qquad&=4\pi\psi0+\int\frac{\mathrm{d}t.\cos q.\psi r}{rr}
\end{aligned}
\]
denotando per $\mathrm{d}t$ indefinite omnia elementa superficiei spatii $s$, atque per $q,r$ eorum respectu eadem, quae antea per literas accentuatas respectu elementorum determinatorum expressa sunt, denique per $\pi$ semicircumferentiam circuli pro radio $=1$.

Ceterum facile perspicitur, si punctum $\mu$ esset neque extra spatium $s$ neque intra, sed in ipsa eius superficie, valere formulam secundam, mutato factore $4\pi$ in $2\pi$, siquidem superficies in puncto $\mu$ neque cuspidem neque aciem offerat; sed ad propositum nostrum haud necessarium est, ad hunc casum attendere.

\Article{8}
Per disquisitionem art. praec. evolutio integralis $\int\int\mathrm{d}s.\mathrm{d}S.\varphi(\mathrm{d}s,\mathrm{d}S)$ reducitur ad
\[
4\pi\sigma\psi0+\int\int\mathrm{d}t.\mathrm{d}S.\frac{\cos q.\psi(\mathrm{d}t,\mathrm{d}S)}{(\mathrm{d}t,\mathrm{d}S)^2}
\]
si per $\sigma$ denotamus volumen eius spatii, quod utrique spatio $s,S$ commune est, ita ut prior pars $4\pi\sigma\psi0$ excidat. si spatia $s,S$ se invicem excludunt. Restat integrale novum, specie etiamnum duplex, revera quintuplex. Quod ut ad quadruplex reducamus, considerabimus integrale
\[
\int\mathrm{d}S.\frac{\cos q.\psi(\mu,\mathrm{d}S)}{(\mu,\mathrm{d}S)^2}
\]
per omnia elementa spatii $S$ extendendum, denotante iterum $\mu$ punctum determinatum, atque $q$ angulum inter duas rectas ab hoc puncto proficiscentes, alteram versus elementum $\mathrm{d}S$, alteram fixam. Hoc integrale, specie simplex, revera triplex, iam ad aliud integrale revera duplex reducere docebimus, et quidem duobus modis prorsus diversis.

Planum per punctum $\mu$ illi rectae fixae normale, per $\Pi$ denotandum, quatenus per proiectionem spatii $S$ attingitur, in elementa infinite parva $\mathrm{d}\Pi$ divisum esse concipiatur. Per punctum talis elementi $\mathrm{d}\Pi$ ducatur recta plano $\Pi$ normalis, quae deinceps, i. e. progrediendo in directione rectae fixae parallela, secet superficiem spatii $S$ in punctis $P',P'',P'''$ etc., quorum distantiae a puncto $\mu$ sint resp. $R',R'',R'''$ etc. Similes rectae per omnia puncta peripheriae ele\SourceLeaf{46}menti $\mathrm{d}\Pi$, plano ad angulos rectos, ductae, spatium prismaticum formabunt, et apud puncta $P',P'',P'''$ etc. e superficie spatii $S$ elementa exsecabunt, quae per $\mathrm{d}T',\mathrm{d}T'',\mathrm{d}T'''$ etc. denotamus. Denique sit $\chi'$ angulus inter duas rectas a puncto $P'$ proficiscentes, alteram extrorsum elemento $\mathrm{d}T'$ normalem, alteram rectae fixae parallelam, similesque angulos apud puncta $P'',P'''$ etc. exprimant characteres $\chi'',\chi'''$ etc. Ita manifesto erit
\[
\mathrm{d}\Pi=-\mathrm{d}T'.\cos\chi'=+\mathrm{d}T''.\cos\chi''=-\mathrm{d}T'''.\cos\chi'''\ \text{etc.}
\]
Spatium prismaticum in elementa infinite parva $\mathrm{d}\Pi.\mathrm{d}z$ dividatur, denotante $z$ distantiam puncti indefiniti a plano $\Pi$ (positive acceptam ab ea parte, a qua est recta fixa); si itaque eiusdem puncti distantiam a puncto $\mu$ per $r$ designamus, erit $z=r\cos q$, nec non (quoniam $rr-zz$ constans est) $z\mathrm{d}z=r\mathrm{d}r$, sive $\mathrm{d}\Pi.\mathrm{d}z.\cos q=\mathrm{d}\Pi.\mathrm{d}r$. Hinc colligitur, integrale nostrum $\int\mathrm{d}S.\frac{\cos q.\psi(\mu,\mathrm{d}S)}{(\mu,\mathrm{d}S)^2}$, extensum per eas spatii $S$ partes, quae in spatio isto prismatico continentur, obtineri per integrale $\mathrm{d}\Pi.\int\frac{\mathrm{d}r.\psi r}{rr}$, si extendatur ab $r=R'$ usque ad $r=R''$, dein ab $r=R'''$ usque ad $r=R''''$ etc. Quodsi itaque indefinite statuimus
\[
\int\frac{\mathrm{d}r.\psi r}{rr}=-\vartheta r
\]
constante integrationis ad lubitum accepta, integrale nostrum pro partibus spatii $S$ intra spatium prismaticum sitis erit
\[
\begin{aligned}
&=\mathrm{d}\Pi.(\vartheta R'-\vartheta R''+\vartheta R'''-\ \text{etc.})\\
&=-\mathrm{d}T'.\cos\chi'.\vartheta R'-\mathrm{d}T''.\cos\chi''.\vartheta R''-\mathrm{d}T'''.\cos\chi'''.\vartheta R'''-\ \text{etc.}
\end{aligned}
\]
Collectis his summationibus per prismata omnibus elementis $\mathrm{d}\Pi$ respondentia, manifesto omnia elementa superficiei spatii $S$ exhausta erunt, habebimusque completum integrale
\[
\int\mathrm{d}S.\frac{\cos q.\psi(\mu,\mathrm{d}S)}{(\mu,\mathrm{d}S)^2}=-\int\mathrm{d}T.\cos\chi.\vartheta R
\]
denotante $\mathrm{d}T$ indefinite quodvis elementum superficiei spatii $S$, $R$ eius distantiam a puncto $\mu$, atque $\chi$ angulum inter normalem ad elementum $\mathrm{d}T$ extrorsum directam atque rectam rectae fixae parallelam.

Hoc itaque modo integrale $\int\int\mathrm{d}s.\mathrm{d}S.\varphi(\mathrm{d}s,\mathrm{d}S)$ reductum est ad formam
\[
4\pi\sigma\psi0-\int\int\mathrm{d}t.\mathrm{d}T.\cos\chi.\vartheta(\mathrm{d}t,\mathrm{d}T)
\]
